\usepackage[a4paper,top=2.5cm,bottom=2.5cm,left=3cm,right=3cm,% margins
    headheight=1.5cm,headsep=1.5em,
    footskip=2em,
]{geometry}

% 载入常用的数学包, 符号包
\usepackage{amsmath}
\usepackage{amsfonts}
\usepackage{amssymb}
\usepackage{mathrsfs}
\usepackage{blindtext}

%% linespace 行间距，段间距等等
\usepackage{setspace}
% \usepackage{indentfirst} % then the first line of each title should start with a indent.
% 定义标题和段落样式
% 定义1.5倍行距
\renewcommand{\baselinestretch}{1.62}
\setlength{\baselineskip}{12pt}   % set the fixed value of the lineskip
\setlength\parskip{\baselineskip} % set the space between the paragraphs, set the variable \parskip \baselineskip
% parindent
\setlength{\parindent}{0pt}

% fonts (style, color, size). 
\usepackage{ctex}		 	% If you are lazy, the CTEX suit is enough.

% Chinese font
\usepackage{xeCJK}		 	        % For the Chinese through XeLaTex
\setCJKmainfont{SimSun} 	        % set the mainfont of Chinese as songti. (serif) for  
\setCJKsansfont{STZhongsong}	            % sans serif font for \textsf
\setCJKmonofont{Maple Font NF CN}	% monospace font for \texttt
% \punctstyle{kaiming}   	        % Remove the space used by symbols like comma. Special for the mainland students like us HUSTers.
\setCJKfamilyfont{song}{SimSun}
\newcommand{\song}{\CJKfamily{song}} %宋体 song
\setCJKfamilyfont{kai}{Kai}
\newcommand{\kai}{\CJKfamily{kai}}   %楷体2312  kai
\setCJKfamilyfont{hwzs}{STZhongsong}
\newcommand{\hwzs}{\CJKfamily{hwzs}} %华文中宋  hwzs
\setCJKfamilyfont{hei}{SimHei}
\newcommand{\hei}{\CJKfamily{hei}}   %黑体  hei

% English font
\usepackage{fontspec}
% font Color 利用definecolor自己可以定义颜色
\usepackage{xcolor}
\definecolor{MSBlue}{rgb}{.204,.353,.541}
\definecolor{MSLightBlue}{rgb}{.31,.506,.741}
% font Size 
\newcommand{\xiaochuhao}{\fontsize{36pt}{\baselineskip}\selectfont}
\newcommand{\erhao}{\fontsize{21pt}{\baselineskip}\selectfont}
\newcommand{\xiaoerhao}{\fontsize{18pt}{\baselineskip}\selectfont}
\newcommand{\sanhao}{\fontsize{15.75pt}{\baselineskip}\selectfont}
\newcommand{\sihao}{\fontsize{14pt}{18pt}\selectfont}
\newcommand{\xiaosihao}{\fontsize{12pt}{18pt}\selectfont}
\newcommand{\wuhao}{\fontsize{10.5pt}{18pt}\selectfont}

%% header and footer 
\usepackage{fancyhdr} % for header and footer
% header
\newcommand{\headstyle}{
    \fancyhead[L]{ \hwzs\wuhao 《class》课程报告}
    \fancyhead[R]{ \hwzs\wuhao id~name}
}
% 设置页脚样式
\newcommand{\footstyle}{\fancyfoot[C]{\wuhao\thepage}
    \fancyfoot[L]{\rule[5pt]{6.7cm}{0.4pt}}
    \fancyfoot[R]{\rule[5pt]{6.7cm}{0.4pt}}
}
\pagestyle{fancy}
\fancyhf{} % 清空原有样式
\headstyle
\footstyle
% 定义一种新的格式叫做main
\fancypagestyle{main}{%
    \fancyhf{} % 清空原有样式
    \headstyle
    \footstyle
}
\renewcommand{\headrulewidth}{0.4pt}
% \renewcommand{\footrulewidth}{0.4pt}
% \renewcommand{\headrule}{\rule{\textwidth}{0.4pt}}

% set the styles of sections at all levels
\usepackage{titlesec}
\usepackage{titletoc}
\titleformat{\section}{\centering\hei\bfseries\xiaoerhao}{\thesection}{1em}{} % 在section标题编号后面加个点
\titleformat*{\subsection}{\raggedright\hei\bfseries\sihao}
\titleformat*{\subsubsection}{\raggedright\hei\bfseries\xiaosihao}
\titleformat{\paragraph}[hang]{\raggedright\hei\bfseries\xiaosihao}{\theparagraph}{1em}{}[]

% manual
% \titleformat{command}[shape]{format}{label}{sep}{before-code}[after-code]
% \titlespacing{command}{left}{before-sep}{after-sep}
% 设置新的层级subsubsubsection
\setcounter{tocdepth}{4}
\setcounter{secnumdepth}{4}

\newcommand{\sectionbreak}{\clearpage} % 小节从新的一页开始
% 根据学校要求设置新的section, subsection, subsection,  paragraph

% set the content of section and so on
\newcommand\seccontent{
    \song
    \xiaosihao % 默认五号字体, 行间距为1.5*\baselineskip
    \setlength{\parindent}{2em} % 首段缩进两个M字符
    \setlength{\parskip}{0pt}
}
\newcommand\tabcontent{
    \song
    \wuhao % 默认五号字体, 行间距为1.5*\baselineskip
    \setlength{\parindent}{2em} % 首段缩进两个M字符
    \setlength{\parskip}{0pt}
}

% ---------------------------------------------------------------------------- %
% for the style of theorems, definitions, proofs and remarks 定义数学里面一些常用的环境
\usepackage{amsthm}
\newtheorem{thm}{\textbf{定理}}[section]
% The section in [] can be replaced by chapter or subsection
\theoremstyle{definition}
\theoremstyle{plain}
\theoremstyle{remark}

% ---------------------------------------------------------------------------- %
% for the caption and reference 图表及公式的编号规范
\usepackage{float} 		 		  	% table figure positioning
\usepackage{caption}
\captionsetup[figure]{labelformat=default, labelsep=quad,name={图}}
\captionsetup[table]{labelformat=default,labelsep=quad,name={表}}
% 设置图表标题的计数方式
\renewcommand{\thefigure}{\ifnum \thesection>0 \thesection-\fi \arabic{figure}} % set caption label style to 2-1
\renewcommand{\thetable}{\ifnum \thesection>0 \thesection-\fi \arabic{table}} % set caption label style to 2-1
\DeclareCaptionFont{mylabelfont}{\hei\xiaosihao}
\captionsetup[figure]{font=mylabelfont}
\captionsetup[table]{font=mylabelfont}

% 设置图表的autoref的格式
\newcommand{\reffig}[1]{图 \ref{#1}}
\newcommand{\reftab}[1]{表 \ref{#1}}
% 公式的编号格式
\renewcommand\theequation{\arabic{section}-\arabic{equation}}

\usepackage{graphicx} % To include graphixs 添加图所需的包
\usepackage{booktabs} % To create three line table including the commands toprule, bottomrule, and midrule
% 使用tabularx库并定义新的左右中格式
\usepackage{tabularx}
\usepackage{makecell}
\newcolumntype{L}{X}
\newcolumntype{C}{>{\centering \arraybackslash}X}
\newcolumntype{R}{>{\raggedright \arraybackslash}X}

% ---------------------------------------------------------------------------- %
% set the style of counters
\makeatletter
\@addtoreset{footnote}{page}
\@addtoreset{figure}{section}
\@addtoreset{table}{section}
\@addtoreset{equation}{section}
\makeatother

% ---------------------------------------------------------------------------- %
% tableofcontents, listoftables and listoffigures 目录
%\renewcommand\listfigurename{插图列表}
%\renewcommand\listtablename{表格列表}
%\titlecontents{标题名}[左间距]{标题格式}{标题标志}{无序号标题}{指引线与页码}[下间距]
%\dottedcontents{section}[2.55em]{\song \xiaosihao \bfseries}{2.5em}{1em}
\usepackage{tocloft}
\renewcommand{\contentsname}{\centerline{ \hei\bfseries\xiaoerhao 目\hspace{2em}录}}
\titlecontents{section}[3em]{\song\xiaosihao\bfseries}{\contentslabel{3em}}{\hspace*{-3em}}{\normalfont\titlerule*[8pt]{.}\contentspage}
\titlecontents{subsection}[3em]{\song\xiaosihao}{\contentslabel{3em}}{\hspace*{-3em}}{\titlerule*[8pt]{.}\contentspage}
\titlecontents{subsubsection}[4em]{\song\xiaosihao}{\contentslabel{4em}}{\hspace*{-4em}}{\titlerule*[8pt]{.}\contentspage}
\titlecontents{paragraph}[5em]{\song\xiaosihao}{\contentslabel{5em}}{\hspace*{-5em}}{\titlerule*[8pt]{.}\contentspage}

% reference and citation 参考文献
\usepackage{natbib}
\renewcommand{\refname}{\centering\hei\xiaoerhao 参考文献}
\bibsep=0pt % 用来设置每个\bibitem之间的间距
% \newcommand{\upcite}[1]{\textsuperscript{\textsuperscript{\cite{#1}}}} % show citation label in the upperscript

% 使用特殊符号
\usepackage{amssymb}
\usepackage{wasysym}

% 制作tatement中的符号
\def\HUSTcheckedbox{$\Square\!\!\!\!\checkmark$}
\def\HUSTbox{$\Square$}

\newcommand{\makestatement}[2]{\statement{#1}{#2}}

\usepackage{enumitem}
\setlist{noitemsep}

\usepackage{listings} % For the code. 代码
\usepackage[bookmarks=true,colorlinks,linkcolor=black,citecolor=black,urlcolor=purple]{hyperref}
\usepackage{appendix}
\renewcommand{\appendixname}{附录}

% titlepage
\usepackage{titling}
% 重置命令 maketitle
\renewcommand{\maketitle}{
    \def\HUSTtitlelength{12em}
    \begin{titlepage}
        \begin{center}
            \vspace*{0em}
            \includegraphics[height=1.61cm]{fig/name.png}\\
            \vspace*{4em}
            {\xiaochuhao \hwzs \bfseries 《class》课程报告}\\
            \vspace*{6em}
            {\erhao \kai \bfseries \thetitle}

            \vspace*{6em}
            {\sanhao \hwzs
                \renewcommand\arraystretch{2}
                \begin{tabular}{lc}
                    % \makebox[4em][s]{院 \hfill 系} &
                    % \underline{\makebox[\HUSTtitlelength]{\school}} \\
                    \makebox[4em][s]{专业}         &
                    \underline{\makebox[\HUSTtitlelength]{\major}}     \\
                    \makebox[4em][s]{班级}         &
                    \underline{\makebox[\HUSTtitlelength]{\class}}     \\
                    \makebox[4em][s]{姓 \hfill 名} &
                    \underline{\makebox[\HUSTtitlelength]{\theauthor}} \\
                    \makebox[4em][s]{学 \hfill 号} &
                    \underline{\makebox[\HUSTtitlelength]{\studentId}} \\
                    % \makebox[4em][s]{指导教师} &
                    % \underline{\makebox[\HUSTtitlelength]{\instructor}} \\
                    % 	\makebox[4em][s]{日\hfill 期} &
                    % \underline{\makebox[\HUSTtitlelength]{\thedate}} \\
                \end{tabular}
            }

            \vspace{4em}
            {\sanhao \hwzs \thedate}
        \end{center}
    \end{titlepage}
}

\newcommand{\hongzifuzhu}[1]{\textcolor{red}{\kai \wuhao(#1)}}